\section{结论}
\label{sec:conclusion}

本文系统阐述了纵列双发正交单轴矢量推力飞行器构型的理论基础，按从推进系统物理到
整机控制律的逻辑链展开。主要工作包括：

\begin{enumerate}[label=(\arabic*)]
    \item \textbf{推进系统建模}：从桨盘动量理论和叶素理论出发，建立了推力和反扭矩的
          集总参数模型（$T = k_T\omega^2$，$\tau = k_Q\omega^2 + J_p\dot{\omega}$），
          阐述了量纲分析的相似律（$k_T \propto \rho D^4$，$k_Q \propto \rho D^5$），
          并分析了电机一阶动态和旋转部件陀螺效应的物理机制；
    \item \textbf{推力矢量映射}：通过摆座运动学分析，导出了完整的六维力/力矩映射
          （式\ref{eq:six_component}），分离了主控通道与耦合通道，
          揭示了推力线过质心和正交摆轴是确保解耦性的必要几何条件；
    \item \textbf{六自由度动力学}：基于牛顿-欧拉方程和四元数姿态运动学，
          建立了整机刚体动力学模型，给出了子步调度的数值积分策略；
    \item \textbf{配平与数值验证}：联立纵向力、法向力与俯仰矩求得24\,m/s纵向工作点
          $\alpha_0=1.627^\circ$、油门$0.441350$、$\delta_{t0}=-0.650^\circ$，
          并由Web core的46项自动测试和三条回归基线验证实现一致性；该点仍有小幅横航向残差，
          不应表述为六自由度全配平；
    \item \textbf{控制效能分析}：给出了含反扭矩交叉项的完整一阶效能矩阵
          （式\ref{eq:effectiveness_matrix}），并将正交摆轴形成的对角主通道近似与完整模型明确区分；
    \item \textbf{稳定性分析框架}：给出了小扰动线性化方法，解释了纵向两模态（短周期、
          长周期）和横航向三模态（滚转收敛、荷兰滚、螺旋）的物理本质，
          以及SAS闭环对模态特性的预期影响，同时明确具体极点仍需数值Jacobian验证；
    \item \textbf{控制律设计}：给出了同时满足控制效率与显示角运动学的三通道SAS增稳律、
          角速度闭环P控制律，以及执行器限幅和积分状态钳位策略。
\end{enumerate}

该构型的核心创新在于：\textbf{通过几何设计（正交摆轴、推力线过质心、反向旋转对消）
形成了结构清晰的主控制通道}，使低复杂度直接反馈成为可行的概念方案。
这一设计哲学——用几何换取算法简洁性——在小型电推进飞行器的工程约束下具有实际价值：
其潜在收益是降低名义工作点的控制分配复杂度；是否无需在线优化、能否实现故障隔离，
仍取决于交叉项补偿、饱和管理和失效状态验证。

本文的解析框架不限定具体型号，但数值实现采用\texttt{models/aircraft-model.json}中的
MODEL-DEFAULT参数，全部未经台架或飞行标定，结论应限定为概念级。
文中导出的解析公式（式\ref{eq:six_component}的六维映射、式\ref{eq:effectiveness_matrix}
的效能矩阵、式\ref{eq:coupling_ratio}的耦合比）可作为参数优化、控制律设计
和子系统选型的理论框架。对于具体型号的工程实现，还需在参数标定（台架试验）、
气动数据获取（CFD或风洞）、实时仿真验证（HIL）和飞行试验等多个维度持续迭代。
在控制律方面，第\ref{sec:INDI}节提出的基于增量非线性动态逆（INDI）的内环升级方案，
利用传感器实测角加速度替代了完整气动模型的依赖，仅需第\ref{sec:allocation}节已推导的
控制效能矩阵$\mathbf{B}_{\mathrm{full}}$，有望在不改变外环SAS结构的前提下显著提升
闭环系统对未建模动态（陀螺力矩、Coriolis耦合、反扭矩交叉项等）的鲁棒性。
该方案的数值仿真验证将作为下一步工作。
